% Zone „Das kennst du schon“ – aus bank/zuordnungen/zone.jsonl (zusammenbau v0.1)
\einheitenkopf[zone][Kennst du schon]{Das kennst du schon}
\setcounter{aufgabe}{0}

%% TODO Titel: Ich-kann-Satz fehlt in der Bank (Platzhalter: Kettenname) – Nr. 1
%% TODO Anweisung: ein Satz über der ersten Teilaufgabe fehlt in der Bank – Nr. 1
\begin{aufgabe}{Punkte im Koordinatensystem eintragen und ablesen (erster Quadrant)}
\begin{teile}
\teil Trage den Punkt A(4|3) ins Koordinatensystem ein.

\begin{ksys}[xmin=0,xmax=6,ymin=0,ymax=6]\end{ksys}
\teil Lies ab: Welche Koordinaten hat der Punkt C?

\begin{ksys}[xmin=0,xmax=6,ymin=0,ymax=6,ablesen]\punkt{2.5}{4}{C}\end{ksys}
C( \leerfeld | \leerfeld )
\end{teile}
\end{aufgabe}

%% TODO Titel: Ich-kann-Satz fehlt in der Bank (Platzhalter: Kettenname) – Nr. 2
%% TODO Anweisung: ein Satz über der ersten Teilaufgabe fehlt in der Bank – Nr. 2
\begin{aufgabe}{Multiplizieren und Dividieren mit Dezimalzahlen (Preis geteilt durch Stückzahl, Preis mal Stückzahl)}
\begin{teile}
\teil 3 Hefte kosten je 1,50 €. Was kosten sie zusammen? \leerfeld[€]
\teil 7 Kugeln Eis kosten je 1,20 €. Was kosten alle zusammen? \leerfeld[€]
\end{teile}
\end{aufgabe}

%% TODO Titel: Ich-kann-Satz fehlt in der Bank (Platzhalter: Kettenname) – Nr. 3
%% TODO Anweisung: ein Satz über der ersten Teilaufgabe fehlt in der Bank – Nr. 3
\begin{aufgabe}{Vielfache und Teiler erkennen (zwölf ist das Dreifache von vier)}
\begin{teile}
\teil 15 ist das Wievielfache von 5? das \leerfeld -Fache
\teil 3 kg sind das Wievielfache von 0,5 kg? das \leerfeld -Fache
\end{teile}
\end{aufgabe}

%% TODO Titel: Ich-kann-Satz fehlt in der Bank (Platzhalter: Kettenname) – Nr. 4
%% TODO Anweisung: ein Satz über der ersten Teilaufgabe fehlt in der Bank – Nr. 4
\begin{aufgabe}{Bruchteil einer Größe (die Hälfte, ein Viertel von 12 €)}
\begin{teile}
\teil Die Hälfte von 14 € – wie viel? \leerfeld[€]
\teil Ein Viertel von 3 € – wie viel? \leerfeld[€]
\end{teile}
\end{aufgabe}

%% TODO Titel: Ich-kann-Satz fehlt in der Bank (Platzhalter: Kettenname) – Nr. 5
%% TODO Anweisung: ein Satz über der ersten Teilaufgabe fehlt in der Bank – Nr. 5
\begin{aufgabe}{Einheiten umrechnen (Cent ↔ Euro, Minuten ↔ Stunden, m ↔ km)}
\begin{teile}
\teil Wie viele Minuten sind 3 Stunden? \leerfeld[min]
\teil Wie viele Meter sind 4,2 km? \leerfeld[m]
\end{teile}
\end{aufgabe}

%% TODO Titel: Ich-kann-Satz fehlt in der Bank (Platzhalter: Kettenname) – Nr. 6
%% TODO Anweisung: ein Satz über der ersten Teilaufgabe fehlt in der Bank – Nr. 6
\begin{aufgabe}{Vielfache und Teiler erkennen (zwölf ist das Dreifache von vier) – Fehler finden}
\begin{teile}
\teil Finn schreibt: „Von 7 € auf 28 € sind es 21 € mehr, also sind 28 € das 21-Fache von 7 €.“ Finde den Fehler und rechne richtig.
\end{teile}
\end{aufgabe}

%% TODO Titel: Ich-kann-Satz fehlt in der Bank (Platzhalter: Kettenname) – Nr. 7
%% TODO Anweisung: ein Satz über der ersten Teilaufgabe fehlt in der Bank – Nr. 7
\begin{aufgabe}{Vielfache und Teiler erkennen (zwölf ist das Dreifache von vier) – selbst rechnen}
\begin{teile}
\teil Ein Beet war 9 m² groß, jetzt ist es 45 m² groß. Das Wievielfache ist das? das \leerfeld -Fache
\end{teile}
\end{aufgabe}
